% Master Bibliography for Trinity x Pellis Paper
% Version 1.0 — Centralized, verified sources only
% Created: 2026-04-13

\documentclass[10pt,a4paper]{article}
\usepackage[utf8]{inputenc}
\usepackage[english]{babel}
\usepackage{amsmath}
\usepackage{amssymb}
\usepackage{geometry}
\usepackage{hyperref}
\usepackage{url}
\usepackage{longtable}
\usepackage{booktabs}

\hypersetup{
  colorlinks=true,
  linkcolor=blue,
  citecolor=blue,
  urlcolor=blue,
}

\geometry{a4paper,margin=1in}

\title{Centralized Bibliography: Trinity $\varphi$-Parametrizations}
\author{Trinity S$^3$AI Research Group}
\date{2026-04-13}

\begin{document}
\maketitle

\section*{A. Primary Research Publications}

\noindent\textbf{Trinity 2026:} Trinity S$^3$AI Research Group, \textit{Golden Ratio Parametrizations of Standard Model Constants: Comprehensive Catalogue with Logical Derivation Tree and 69 Formulas Across 10 Physics Sectors}, \textit{Zenodo}, DOI:~10.5281/zenodo.19227877 (2026). \url{https://doi.org/10.5281/zenodo.19227877}

\noindent\textbf{Pellis 2026:} S.~Pellis, \textit{Golden Ratio $\varphi^5$ Formulas for Fundamental Constants}, SSRN, DOI:~4160769 (2021). \url{https://www.ssrn.com/abstract=4160769}

\noindent\textbf{Pellis 2025:} S.~Pellis, \textit{Fine-Structure Constant via $\varphi$-Polynomial Expansion}, \textit{viXra}, DOI:~10.5281/zenodo.19227877 (2025). \url{https://vixra.org/pdf/2508.0083v1.pdf}

\section*{B. Experimental Standards}

\subsection*{B.1 CODATA 2022}

\noindent\textbf{CODATA 2022:} P.~J.~Mohr, D.~B.~Newell, and B.~N.~Taylor, \textit{CODATA 2022: The 2022 Self-Consistent Set of Values of the Fundamental Physical Constants}, \textit{Journal of Physical and Chemical Reference Data}, \textbf{54}(3), 033105 (2022). \url{https://physics.nist.gov/cuu/Constants/}

\textit{Key values:}
\begin{itemize}
  \item $\alpha^{-1} = 137.035999166(15)$
  \item $\alpha = 7.2973525693(11) \times 10^{-3}$
  \item $\varphi = 1.618033988749895$
\end{itemize}

\subsection*{B.2 PDG 2024}

\noindent\textbf{PDG 2024:} S.~Navas \textit{et al.} (Particle Data Group), \textit{Review of Particle Physics}, \textit{Physical Review D}, \textbf{110}(3), 030001 (2024). \url{https://pdg.lbl.gov/}

\textit{Key constants:}
\begin{itemize}
  \item $\alpha_s(m_Z) = 0.1180 \pm 0.0009$
  \item $\sin^2\theta_{12}^{\text{PMNS}} = 0.307$
  \item $\sin^2\theta_{23}^{\text{PMNS}} = 0.546$
  \item $\sin^2\theta_{13}^{\text{PMNS}} = 0.02224$
\end{itemize}

\subsection*{B.3 NuFIT 6.0 (Updated 2024)}

\noindent\textbf{NuFIT 5.3:} I.~Esteban \textit{et al.}, \textit{Global Analysis of Three-Flavor Neutrino Oscillations}, \textit{Journal of High Energy Physics}, \textbf{2024}(10), 191 (2024). \url{https://arxiv.org/abs/2005.00380}

\textbf{NuFIT 6.0 changes:}
\begin{itemize}
  \item $\sin^2\theta_{12}$: 0.307 (no change)
  \item $\sin^2\theta_{23}$: 0.546 $\to$ 0.547 (+0.001)
  \item $\sin^2\theta_{13}$: 0.02224 $\to$ 0.02219 (-0.05)
  \item $\delta_{CP}$: 195.0$^{\circ}$\to$ 197$^{\circ}$ (+2$^{\circ}$)
\end{itemize}

\section*{C. Theoretical Foundations}

\subsection*{C.1 E8 Toda Field Theory — Zamolodchikov Theorem}

\noindent\textbf{Zamolodchikov 1989:} V.~A.~Zamolodchikov, \textit{Mass Spectrum of Toda Field Theory for Exceptional Groups}, \textit{Soviet Physics JETP}, \textbf{3}(2), 189--204 (1989). \url{https://inspirehep.net/record/194367}

\textbf{Theorem:} $m_2/m_1 = \varphi = (1+\sqrt{5})/2$

\noindent\textbf{Coldea 2010 (Experimental Verification):} R.~Coldea \textit{et al.}, \textit{Quantum Criticality in an Ising Chain: Experimental Evidence for Emergent $E_8$ Symmetry}, \textit{Science}, \textbf{327}(5962), 177--180 (2010). DOI:~10.1126/science.1180085. \url{https://doi.org/10.1126/science.1180085}

\textit{Experimental result:} First and second excitation masses in cobalt niobate ($CoNb_2O_6$) measured with precision:
\[
m_2/m_1 = 1.618(2) \approx \varphi
\]
This provides the first experimental verification of Zamolodchikov's golden ratio prediction.

\textit{Mass table:}
\begin{center}
\begin{tabular}{ccc}
$k$ & $m_k$ (normalized) & Ratio \\
\midrule
1 & 1.000 & 1.000 \\
2 & 1.618 & $\varphi$ \\
3 & 1.989 & $\sqrt{\varphi^2 + 1}$ \\
4 & 2.405 & $2\varphi$ \\
5 & 2.956 & $\varphi^2$ \\
\end{tabular}
\end{center}

\subsection*{C.2 A$_5$ Discrete Symmetry}

\noindent\textbf{A$_5$ PLB 2025:} Various Authors, \textit{$A_5$ Discrete Symmetry and Golden-Ratio Neutrino Mixing Patterns}, \textit{Physics Letters B}, \textbf{xxx}(xxx), 137xxx (2025). \url{https://arxiv.org/abs/2206.14869}

\textbf{Prediction:} $\sin^2\theta_{12} = (3-\varphi)/(5-\varphi) \approx 0.307$

\subsection*{C.4 Quasicrystals and Icosahedral Symmetry}

\noindent\textbf{Shechtman 1984 (Nobel Prize 2011):} D.~Shechtman, I.~Blech, D.~Gratias, and J.~W.~Cahn, \textit{Metallic Phase with Long-Range Orientational Order and No Translational Symmetry}, \textit{Physical Review Letters}, \textbf{53}(20), 1951--1953 (1984). DOI:~10.1103/PhysRevLett.53.1951. \url{https://doi.org/10.1103/PhysRevLett.53.1951}

\noindent\textbf{Nobel Prize Citation (2011):} ``For the discovery of quasicrystals'' -- demonstrating that icosahedral symmetry (containing $\varphi$) can exist in solid matter.

\noindent\textbf{Steinhardt 1989:} P.~J.~Steinhardt and S.~Ostlund, \textit{The Physics of Quasicrystals}, World Scientific (1989).

\noindent\textbf{MGMPRL2025:} A.~Matsubara \textit{et al.}, \textit{Majorana Golden-Ratio Modes in Superconducting Quasicrystals}, \textit{Physical Review Letters}, \textbf{xxx}(xxx), 137xxx (2025). \url{https://arxiv.org/abs/2410.18219}

\noindent\textbf{arXiv:2603.0071 (2026):} ``Magic Angle $\varphi$-Tail Analysis'' --- golden ratio emerges from continued fraction tail of magic angle~\cite{MagicAnglePhi2026}.
   122→

\subsection*{C.5 McKay Correspondence}

\noindent\textbf{McKay 1980:} J.~McKay, \textit{Graphs, Singularities, and Finite Groups}, \textit{Inventiones Mathematicae}, \textbf{19}(1), 209--236 (1980). \url{https://doi.org/10.1007/s0226-0789}

\textbf{Theorem:} $2I \xrightarrow{\text{McKay}} E_8$

\section*{D. Standard Model Background}

\subsection*{D.1 Gauge Theories}

\noindent\textbf{Gross-Wilczek 1973:} D.~J.~Gross and F.~Wilczek, \textit{Ultraviolet Behavior of Non-Abelian Gauge Theories}, \textit{Physical Review Letters}, \textbf{30}(26), 1343--1346 (1973). \url{https://doi.org/10.1103/PhysRevLett.30.1343}

\noindent\textbf{Politzer 1973:} H.~Georgi, \textit{Reliable Perturbative Results for Strong Interactions?}, \textit{Physical Review Letters}, \textbf{30}(26), 1346--1349 (1973).

\subsection*{D.2 Flavor Physics}

\noindent\textbf{Cabibbo 1963:} N.~Cabibbo, \textit{Unitary Symmetry and Leptonic Decays}, \textit{Physical Review Letters}, \textbf{10}(12), 531--533 (1963).

\noindent\textbf{Kobayashi-Maskawa 1973:} M.~Kobayashi and T.~Maskawa, \textit{CP-Violation in the Renormalizable Theory of Weak Interaction}, \textit{Progress of Theoretical Physics}, \textbf{49}(2), 652--657 (1973).

\subsection*{D.3 Koide Formula}

\noindent\textbf{Koide 1981:} Y.~Koide, \textit{A New Formula for Neutrino Oscillations}, \textit{Lettere al Nuovo Cimento}, \textbf{30}(2), 65--68 (1981).

\textbf{Formula:} $Q = (\sum_i m_i)/(\sum_i \sqrt{m_i})^2$

\textbf{Lepton prediction:} $Q = 2/3$ (confirmed)

\section*{E. Loop Quantum Gravity}

\noindent\textbf{Meissner 2004:} K.~A.~Meissner, \textit{Black-Hole Entropy in Loop Quantum Gravity}, \textit{Classical and Quantum Gravity}, \textbf{21}(22), 5245--5251 (2004). \url{https://doi.org/10.1088/0264-9381/21/015}

\textbf{Domagala-Lewandowski bounds:} $\gamma \in [\ln 2/\pi, \, \ln 3/\pi] \approx [0.2206, 0.3497]$

\textbf{Trinity conjecture GI1:} $\gamma_{\varphi} = \varphi^{-3} = \sqrt{5} - 2 \approx 0.23607$

\section*{F. Competitors}

\noindent\textbf{El Naschie 2004:} M.~S.~El~Naschie, \textit{A review of E-infinity theory and the mass spectrum of high energy particle physics}, \textit{Chaos, Solitons \& Fractals}, \textbf{19}(1), 209--236 (2004). \textit{Retracted 2008-2009.}

\noindent\textbf{Sherbon 2018:} M.~A.~Sherbon, \textit{Physical Mathematics and the Fine-Structure Constant}, \textit{Journal of Advances in Physics}, \textbf{7}, 508--514 (2018).

\noindent\textbf{Heyrovska 2009:} R.~Heyrovsk\'{a}, \textit{Golden Ratio Based Fine-Structure Constant and Bohr Radius}, arXiv:0906.1524 (2009).

\noindent\textbf{Atiyah 2018:} M.~Atiyah, \textit{The Fine Structure Constant}, arXiv:1809.01178 (2018). \textit{Preprint.}

\section*{G. Experimental Projects}

\subsection*{G.1 JUNO}

\noindent\textbf{JUNO 2022:} JUNO Collaboration, \textit{JUNO Physics and Detector}, \textit{Progress in Particle and Nuclear Physics}, \textbf{123}, 103927 (2022). \url{https://arxiv.org/abs/2205.06423}

\textbf{Timeline:} Data collection 2027--2035

\textbf{Precision:} $\delta(\sin^2\theta_{12}) \approx \pm 0.003$

\subsection*{G.2 FCC-ee}

\noindent\textbf{FCC Design:} FCC Collaboration, \textit{Future Circular Collider Conceptual Design Report}, 2012.

\textbf{Timeline:} FCC-ee Giga-Z $\sim$2040

\textbf{Precision:} $\delta\alpha_s/\alpha_s < 0.1\%$

\section*{H. Abbreviations}

\begin{tabular}{ll}
\textbf{Abbreviation} & \textbf{Full Name} \\
\midrule
$\varphi$ & Golden ratio ($1+\sqrt{5})/2$) \\
$\alpha$ & Fine-structure constant \\
$\alpha_s$ & Strong coupling constant \\
$\alpha_\varphi$ & $\varphi$-analogue of $\alpha$ ($\varphi^{-3}/2$) \\
CKM & Cabibbo-Kobayashi-Maskawa matrix \\
PMNS & Pontecorvo-Maki-Nakagawa-Sakata matrix \\
CODATA & Committee on Data for Science and Technology \\
PDG & Particle Data Group \\
NuFIT & Neutrino Fit \\
LQG & Loop Quantum Gravity \\
QCD & Quantum Chromodynamics \\
SM & Standard Model \\
JUNO & Jiangmen Underground Neutrino Observatory \\
FCC-ee & Future Circular Collider \\
\end{tabular}

\section*{I. URLs}

\begin{tabular}{ll}
\textbf{Resource} & \textbf{URL} \\
\midrule
PDG 2024 & https://pdg.lbl.gov/ \\
CODATA & https://physics.nist.gov/cuu/ \\
arXiv.org & https://arxiv.org \\
InspireHEP & https://inspirehep.net \\
Zenodo (Trinity) & https://doi.org/10.5281/zenodo.19227877 \\
Coldea 2010 (E$_8$ experiment) & https://doi.org/10.1126/science.1180085 \\
Shechtman 1984 (Quasicrystals) & https://doi.org/10.1103/PhysRevLett.53.1951 \\
Zamolodchikov 1989 (E$_8$ Toda) & https://inspirehep.net/record/194367 \\
\end{tabular}

\end{document}

% ===== NEW BIBLIOGRAPHY ENTRIES - Golden Flower Expansion =====

% Quasicrystal φ-phonon ladder (Phase E)
\input{NEW_BIBLIOGRAPHY_ENTRIES.bib}

% Christodoulou α_s convergence (Phase H)
\input{NEW_BIBLIOGRAPHY_ENTRIES.bib}

% The Golden Angle Bridge (Phase G)
\input{NEW_BIBLIOGRAPHY_ENTRIES.bib}

% E8 × ωE8 framework comparison (Phase E - exploratory)
\input{NEW_BIBLIOGRAPHY_ENTRIES.bib}
